% !TeX program = lualatex
% !TeX root = predavanje05.tex
% ============================================================
% Zajednički preamble za sva predavanja RWA kolegija
% Uključiti s: % ============================================================
% Zajednički preamble za sva predavanja RWA kolegija
% Uključiti s: % ============================================================
% Zajednički preamble za sva predavanja RWA kolegija
% Uključiti s: \input{../common/preamble.tex}
% ============================================================

\documentclass[aspectratio=149,11pt]{beamer}
\usetheme{CambridgeUS}

\usepackage{fontspec}              % radi samo sa XeLaTeX/LuaLaTeX
\usepackage[croatian]{babel}
\usepackage{csquotes}              % \enquote{} → hrvatski navodnici „..."
\usepackage{booktabs}
\usepackage{array}
\usepackage{listings}
\usepackage{xcolor}
\usepackage{colortbl}
\usepackage{hyperref}
\usepackage{tikz}
\usetikzlibrary{positioning,arrows.meta}

% ---------- Tipografija: IBM Plex ----------
\setsansfont{IBMPlexSans}[
  Extension      = .otf,
  UprightFont    = *-Regular,
  ItalicFont     = *-Italic,
  BoldFont       = *-SemiBold,
  BoldItalicFont = *-SemiBoldItalic
]
\setmainfont{IBMPlexSans}[
  Extension      = .otf,
  UprightFont    = *-Regular,
  ItalicFont     = *-Italic,
  BoldFont       = *-SemiBold,
  BoldItalicFont = *-SemiBoldItalic
]
\setmonofont{IBMPlexMono}[
  Extension      = .otf,
  UprightFont    = *-Regular,
  ItalicFont     = *-Italic,
  BoldFont       = *-SemiBold,
  BoldItalicFont = *-SemiBoldItalic,
  Scale          = 0.92
]

% ---------- Paleta ----------
\definecolor{ink}{HTML}{0C4685}      % primarna tamnoplava
\definecolor{ink2}{HTML}{6E7BAA}     % sekundarna plavo-siva
\definecolor{tint}{HTML}{EAEEF6}     % svjetli plavi tint
\definecolor{textclr}{HTML}{1A222E}  % glavni tekst
\definecolor{mutedclr}{HTML}{6B7484} % podnožje, sekundarni tekst
\definecolor{bgclr}{HTML}{F6F7F9}    % pozadina slajda
\definecolor{lineclr}{HTML}{D9DEE7}  % tanke linije tablica
\definecolor{divrule}{HTML}{9DB4D6}  % akcent na tamnim sekcijskim slajdovima
\definecolor{divtext}{HTML}{E6ECF6}  % tekst na tamnim sekcijskim slajdovima

\setbeamercolor{normal text}{fg=textclr,bg=bgclr}
\setbeamercolor{structure}{fg=ink}
\setbeamercolor{alerted text}{fg=ink}

\setbeamercolor{title}{fg=ink,bg=tint}
\setbeamercolor{subtitle}{fg=mutedclr}
\setbeamercolor{author}{fg=textclr}
\setbeamercolor{date}{fg=mutedclr}
\setbeamercolor{frametitle}{fg=ink,bg=tint}

\setbeamercolor{section in head/foot}{fg=white,bg=ink}
\setbeamercolor{subsection in head/foot}{fg=ink,bg=tint}
\setbeamercolor{author in head/foot}{fg=white,bg=ink}
\setbeamercolor{title in head/foot}{fg=textclr,bg=tint}
\setbeamercolor{date in head/foot}{fg=mutedclr,bg=tint}

\setbeamercolor{block title}{fg=white,bg=ink}
\setbeamercolor{block body}{fg=textclr,bg=tint}
\setbeamercolor{block title example}{fg=white,bg=ink2}
\setbeamercolor{block body example}{fg=textclr,bg=tint}
\setbeamercolor{itemize item}{fg=ink2}
\setbeamercolor{itemize subitem}{fg=ink2}
\setbeamercolor{enumerate item}{fg=ink2}

\setbeamertemplate{itemize items}[square]
\setbeamertemplate{navigation symbols}{}
\arrayrulecolor{lineclr}

% ---------- Kolegij, autori i logo ----------
\newcommand{\kolegij}{Programiranje za web}
\author{doc.\,dr.\,sc. Tomislav Petković \and mag.\,ing. Josip Torić}
\newcommand{\znak}[2]{\raisebox{-0.5\height}{\includegraphics[height=#1]{#2}}}
\titlegraphic{\znak{2.1cm}{../common/SIT_logo_tekst1.pdf}\hspace{1.8em}\znak{1.8cm}{../common/UNIZD_logo.pdf}}

% ---------- Zaglavlje ----------
% Jedna puna traka: lijevo sekcija, desno podsekcija.
\setbeamertemplate{headline}{%
  \leavevmode%
  \hbox{%
    \begin{beamercolorbox}[wd=\paperwidth,ht=2.25ex,dp=1ex,leftskip=.3cm,rightskip=.3cm]{section in head/foot}%
      \usebeamerfont{section in head/foot}\insertsectionhead\hfill\insertsubsectionhead%
    \end{beamercolorbox}%
  }%
}

% ---------- Podnožje ----------
% Lijevo: naziv trenutne prezentacije (kratki oblik naslova iz \title[...]{...})
% Sredina: naziv kolegija.  Desno: datum i redni broj slajda.
\setbeamertemplate{footline}{%
  \leavevmode%
  \hbox{%
    \begin{beamercolorbox}[wd=.34\paperwidth,ht=2.25ex,dp=1ex,center]{author in head/foot}%
      \usebeamerfont{author in head/foot}\insertshorttitle%
    \end{beamercolorbox}%
    \begin{beamercolorbox}[wd=.33\paperwidth,ht=2.25ex,dp=1ex,center]{title in head/foot}%
      \usebeamerfont{title in head/foot}\kolegij%
    \end{beamercolorbox}%
    \begin{beamercolorbox}[wd=.33\paperwidth,ht=2.25ex,dp=1ex,center]{date in head/foot}%
      \usebeamerfont{date in head/foot}\insertshortdate{}\hspace*{2em}%
      \insertframenumber{}\,/\,\inserttotalframenumber\hspace*{1ex}%
    \end{beamercolorbox}%
  }%
  \vskip0pt%
}

% ---------- Sekcijski slajd ----------
% \sekcija{Naslov sekcije}{jedna linija koja nabraja teme sekcije}
\newcommand{\sekcija}[2]{%
  \section{#1}%
  {%
    \setbeamercolor{background canvas}{bg=ink}%
    \setbeamertemplate{headline}{}%
    \begin{frame}
      \vfill
      \hspace*{0.03\textwidth}{\color{divrule}\rule{2.6cm}{1.1pt}}\par
      \vskip 0.7em
      \hspace*{0.03\textwidth}{\Large\bfseries\color{white}#1}\par
      \vskip 0.6em
      \hspace*{0.03\textwidth}\begin{minipage}{0.9\textwidth}\color{divtext}\small #2\end{minipage}\par
      \vfill
    \end{frame}%
  }%
}

% ---------- Listings ----------
\lstdefinelanguage{bash}{
  morekeywords={sudo,apt,brew,pip,pipx,uv,winget,python,git,flask,source,deactivate,code,gunicorn,docker,compose,heroku,render,railway,node,npm,pnpm,yarn,uvicorn,alembic,pytest,ruff,black,precommit},
  sensitive=false,
  morecomment=[l]{\#}
}
\lstdefinelanguage{python}{
  morekeywords={from,import,def,class,return,if,elif,else,for,in,while,try,except,raise,with,as,not,and,or,is,lambda,yield,True,False,None,await,async},
  sensitive=true,
  morecomment=[l]{\#},
  morestring=[b]',
  morestring=[b]"
}
\lstdefinelanguage{sql}{
  morekeywords={SELECT,FROM,WHERE,INSERT,INTO,VALUES,UPDATE,SET,DELETE,CREATE,TABLE,PRIMARY,KEY,FOREIGN,REFERENCES,UNIQUE,NOT,NULL,INDEX,ON,ALTER,DROP,CONSTRAINT},
  sensitive=false,
  morecomment=[l]{--}
}
\lstdefinelanguage{json}{
  morecomment=[l]{//},
  morestring=[b]",
  showstringspaces=false
}
\lstdefinelanguage{typescript}{
  morekeywords={import,export,from,const,let,var,function,return,if,else,for,of,in,while,switch,case,break,async,await,try,catch,finally,throw,new,class,this,extends,implements,interface,type,enum,readonly,public,private,protected,static,get,set,true,false,null,undefined,void,string,number,boolean,any,unknown,never,typeof,keyof,as},
  sensitive=true,
  morecomment=[l]{//},
  morecomment=[s]{/*}{*/},
  morestring=[b]',
  morestring=[b]",
  morestring=[b]`
}
\lstalias{javascript}{typescript}
\lstset{
  backgroundcolor=\color{tint},
  basicstyle=\ttfamily\footnotesize,
  keywordstyle=\color{ink}\bfseries,
  commentstyle=\color{mutedclr},
  frame=single,
  rulecolor=\color{lineclr},
  frameround=tttt,
  breaklines=true,
  tabsize=2,
  showstringspaces=false,
  upquote=true,
  columns=fullflexible,
  escapeinside={<@}{@>}
}
% ---------- Dijakritici u isječcima koda ----------
% listings čita bajt po bajt, pa višebajtni znak premjesti: "početni" ispadne
% "čpoetni". literate to ne popravlja. Rješenje je izlaz u LaTeX: sve između
% <@ i @> ne prolazi kroz listings, nego se slaže kao običan tekst, a unutar
% okvira je font ionako ttfamily, pa se ništa ne razlikuje.
%
%   async def get(...):
%       club = await repo.get_by_id(db, club_id)   <@\kom{dohvat kluba}@>
%       raise AppError("not_found", <@"Klub nije pronađen"@>, 404)
%
% Graničnici <@ i @> odabrani su jer se ne pojavljuju ni u jednom jeziku
% koji se na predavanjima citira; (* bi se sudaralo s def f(*args).
\newcommand{\kom}[1]{{\color{mutedclr}\# #1}}
% Isto za TypeScript, Vue i CSS, kojima komentar počinje s //:
%   const lifters = useLiftersStore()   <@\komts{dijeljeni popis}@>
\newcommand{\komts}[1]{{\color{mutedclr}// #1}}

% ---------- Činjenica naspram konvencije ----------
% Sve što nije propisano standardom ili jezikom, nego je dogovor struke,
% ide u ovaj okvir, da student zna koje je vrste tvrdnja pred njim.
% \begin{konvencija} ... \end{konvencija}
% \begin{konvencija}[Konvencija: nazivi putanja] ... \end{konvencija}
\newenvironment{konvencija}[1][Konvencija]
  {\begin{exampleblock}{#1}}
  {\end{exampleblock}}

% ---------- Helper macros ----------
% Stupac tablice zadane širine, poravnat ulijevo (bez rastezanja i dijeljenja riječi).
\newcolumntype{L}[1]{>{\raggedright\arraybackslash}p{#1}}
\newcommand{\codefile}[1]{\texttt{#1}}
\newcommand{\kbd}[1]{\texttt{#1}}

% ---------- Code frame environment ----------
% Automatski dodaje fragile i shrink=5, koristi se umjesto \begin{frame}
% Primjer: \begin{codeframe}{Naslov} ... \end{codeframe}
% Primjer s opcijama: \begin{codeframe}[shrink=15]{Naslov} ... \end{codeframe}
\newenvironment{codeframe}[2][]
  {\begin{frame}[fragile,shrink=5,#1]{#2}}
  {\end{frame}}

% ============================================================

\documentclass[aspectratio=149,11pt]{beamer}
\usetheme{CambridgeUS}

\usepackage{fontspec}              % radi samo sa XeLaTeX/LuaLaTeX
\usepackage[croatian]{babel}
\usepackage{csquotes}              % \enquote{} → hrvatski navodnici „..."
\usepackage{booktabs}
\usepackage{array}
\usepackage{listings}
\usepackage{xcolor}
\usepackage{colortbl}
\usepackage{hyperref}
\usepackage{tikz}
\usetikzlibrary{positioning,arrows.meta}

% ---------- Tipografija: IBM Plex ----------
\setsansfont{IBMPlexSans}[
  Extension      = .otf,
  UprightFont    = *-Regular,
  ItalicFont     = *-Italic,
  BoldFont       = *-SemiBold,
  BoldItalicFont = *-SemiBoldItalic
]
\setmainfont{IBMPlexSans}[
  Extension      = .otf,
  UprightFont    = *-Regular,
  ItalicFont     = *-Italic,
  BoldFont       = *-SemiBold,
  BoldItalicFont = *-SemiBoldItalic
]
\setmonofont{IBMPlexMono}[
  Extension      = .otf,
  UprightFont    = *-Regular,
  ItalicFont     = *-Italic,
  BoldFont       = *-SemiBold,
  BoldItalicFont = *-SemiBoldItalic,
  Scale          = 0.92
]

% ---------- Paleta ----------
\definecolor{ink}{HTML}{0C4685}      % primarna tamnoplava
\definecolor{ink2}{HTML}{6E7BAA}     % sekundarna plavo-siva
\definecolor{tint}{HTML}{EAEEF6}     % svjetli plavi tint
\definecolor{textclr}{HTML}{1A222E}  % glavni tekst
\definecolor{mutedclr}{HTML}{6B7484} % podnožje, sekundarni tekst
\definecolor{bgclr}{HTML}{F6F7F9}    % pozadina slajda
\definecolor{lineclr}{HTML}{D9DEE7}  % tanke linije tablica
\definecolor{divrule}{HTML}{9DB4D6}  % akcent na tamnim sekcijskim slajdovima
\definecolor{divtext}{HTML}{E6ECF6}  % tekst na tamnim sekcijskim slajdovima

\setbeamercolor{normal text}{fg=textclr,bg=bgclr}
\setbeamercolor{structure}{fg=ink}
\setbeamercolor{alerted text}{fg=ink}

\setbeamercolor{title}{fg=ink,bg=tint}
\setbeamercolor{subtitle}{fg=mutedclr}
\setbeamercolor{author}{fg=textclr}
\setbeamercolor{date}{fg=mutedclr}
\setbeamercolor{frametitle}{fg=ink,bg=tint}

\setbeamercolor{section in head/foot}{fg=white,bg=ink}
\setbeamercolor{subsection in head/foot}{fg=ink,bg=tint}
\setbeamercolor{author in head/foot}{fg=white,bg=ink}
\setbeamercolor{title in head/foot}{fg=textclr,bg=tint}
\setbeamercolor{date in head/foot}{fg=mutedclr,bg=tint}

\setbeamercolor{block title}{fg=white,bg=ink}
\setbeamercolor{block body}{fg=textclr,bg=tint}
\setbeamercolor{block title example}{fg=white,bg=ink2}
\setbeamercolor{block body example}{fg=textclr,bg=tint}
\setbeamercolor{itemize item}{fg=ink2}
\setbeamercolor{itemize subitem}{fg=ink2}
\setbeamercolor{enumerate item}{fg=ink2}

\setbeamertemplate{itemize items}[square]
\setbeamertemplate{navigation symbols}{}
\arrayrulecolor{lineclr}

% ---------- Kolegij, autori i logo ----------
\newcommand{\kolegij}{Programiranje za web}
\author{doc.\,dr.\,sc. Tomislav Petković \and mag.\,ing. Josip Torić}
\newcommand{\znak}[2]{\raisebox{-0.5\height}{\includegraphics[height=#1]{#2}}}
\titlegraphic{\znak{2.1cm}{../common/SIT_logo_tekst1.pdf}\hspace{1.8em}\znak{1.8cm}{../common/UNIZD_logo.pdf}}

% ---------- Zaglavlje ----------
% Jedna puna traka: lijevo sekcija, desno podsekcija.
\setbeamertemplate{headline}{%
  \leavevmode%
  \hbox{%
    \begin{beamercolorbox}[wd=\paperwidth,ht=2.25ex,dp=1ex,leftskip=.3cm,rightskip=.3cm]{section in head/foot}%
      \usebeamerfont{section in head/foot}\insertsectionhead\hfill\insertsubsectionhead%
    \end{beamercolorbox}%
  }%
}

% ---------- Podnožje ----------
% Lijevo: naziv trenutne prezentacije (kratki oblik naslova iz \title[...]{...})
% Sredina: naziv kolegija.  Desno: datum i redni broj slajda.
\setbeamertemplate{footline}{%
  \leavevmode%
  \hbox{%
    \begin{beamercolorbox}[wd=.34\paperwidth,ht=2.25ex,dp=1ex,center]{author in head/foot}%
      \usebeamerfont{author in head/foot}\insertshorttitle%
    \end{beamercolorbox}%
    \begin{beamercolorbox}[wd=.33\paperwidth,ht=2.25ex,dp=1ex,center]{title in head/foot}%
      \usebeamerfont{title in head/foot}\kolegij%
    \end{beamercolorbox}%
    \begin{beamercolorbox}[wd=.33\paperwidth,ht=2.25ex,dp=1ex,center]{date in head/foot}%
      \usebeamerfont{date in head/foot}\insertshortdate{}\hspace*{2em}%
      \insertframenumber{}\,/\,\inserttotalframenumber\hspace*{1ex}%
    \end{beamercolorbox}%
  }%
  \vskip0pt%
}

% ---------- Sekcijski slajd ----------
% \sekcija{Naslov sekcije}{jedna linija koja nabraja teme sekcije}
\newcommand{\sekcija}[2]{%
  \section{#1}%
  {%
    \setbeamercolor{background canvas}{bg=ink}%
    \setbeamertemplate{headline}{}%
    \begin{frame}
      \vfill
      \hspace*{0.03\textwidth}{\color{divrule}\rule{2.6cm}{1.1pt}}\par
      \vskip 0.7em
      \hspace*{0.03\textwidth}{\Large\bfseries\color{white}#1}\par
      \vskip 0.6em
      \hspace*{0.03\textwidth}\begin{minipage}{0.9\textwidth}\color{divtext}\small #2\end{minipage}\par
      \vfill
    \end{frame}%
  }%
}

% ---------- Listings ----------
\lstdefinelanguage{bash}{
  morekeywords={sudo,apt,brew,pip,pipx,uv,winget,python,git,flask,source,deactivate,code,gunicorn,docker,compose,heroku,render,railway,node,npm,pnpm,yarn,uvicorn,alembic,pytest,ruff,black,precommit},
  sensitive=false,
  morecomment=[l]{\#}
}
\lstdefinelanguage{python}{
  morekeywords={from,import,def,class,return,if,elif,else,for,in,while,try,except,raise,with,as,not,and,or,is,lambda,yield,True,False,None,await,async},
  sensitive=true,
  morecomment=[l]{\#},
  morestring=[b]',
  morestring=[b]"
}
\lstdefinelanguage{sql}{
  morekeywords={SELECT,FROM,WHERE,INSERT,INTO,VALUES,UPDATE,SET,DELETE,CREATE,TABLE,PRIMARY,KEY,FOREIGN,REFERENCES,UNIQUE,NOT,NULL,INDEX,ON,ALTER,DROP,CONSTRAINT},
  sensitive=false,
  morecomment=[l]{--}
}
\lstdefinelanguage{json}{
  morecomment=[l]{//},
  morestring=[b]",
  showstringspaces=false
}
\lstdefinelanguage{typescript}{
  morekeywords={import,export,from,const,let,var,function,return,if,else,for,of,in,while,switch,case,break,async,await,try,catch,finally,throw,new,class,this,extends,implements,interface,type,enum,readonly,public,private,protected,static,get,set,true,false,null,undefined,void,string,number,boolean,any,unknown,never,typeof,keyof,as},
  sensitive=true,
  morecomment=[l]{//},
  morecomment=[s]{/*}{*/},
  morestring=[b]',
  morestring=[b]",
  morestring=[b]`
}
\lstalias{javascript}{typescript}
\lstset{
  backgroundcolor=\color{tint},
  basicstyle=\ttfamily\footnotesize,
  keywordstyle=\color{ink}\bfseries,
  commentstyle=\color{mutedclr},
  frame=single,
  rulecolor=\color{lineclr},
  frameround=tttt,
  breaklines=true,
  tabsize=2,
  showstringspaces=false,
  upquote=true,
  columns=fullflexible,
  escapeinside={<@}{@>}
}
% ---------- Dijakritici u isječcima koda ----------
% listings čita bajt po bajt, pa višebajtni znak premjesti: "početni" ispadne
% "čpoetni". literate to ne popravlja. Rješenje je izlaz u LaTeX: sve između
% <@ i @> ne prolazi kroz listings, nego se slaže kao običan tekst, a unutar
% okvira je font ionako ttfamily, pa se ništa ne razlikuje.
%
%   async def get(...):
%       club = await repo.get_by_id(db, club_id)   <@\kom{dohvat kluba}@>
%       raise AppError("not_found", <@"Klub nije pronađen"@>, 404)
%
% Graničnici <@ i @> odabrani su jer se ne pojavljuju ni u jednom jeziku
% koji se na predavanjima citira; (* bi se sudaralo s def f(*args).
\newcommand{\kom}[1]{{\color{mutedclr}\# #1}}
% Isto za TypeScript, Vue i CSS, kojima komentar počinje s //:
%   const lifters = useLiftersStore()   <@\komts{dijeljeni popis}@>
\newcommand{\komts}[1]{{\color{mutedclr}// #1}}

% ---------- Činjenica naspram konvencije ----------
% Sve što nije propisano standardom ili jezikom, nego je dogovor struke,
% ide u ovaj okvir, da student zna koje je vrste tvrdnja pred njim.
% \begin{konvencija} ... \end{konvencija}
% \begin{konvencija}[Konvencija: nazivi putanja] ... \end{konvencija}
\newenvironment{konvencija}[1][Konvencija]
  {\begin{exampleblock}{#1}}
  {\end{exampleblock}}

% ---------- Helper macros ----------
% Stupac tablice zadane širine, poravnat ulijevo (bez rastezanja i dijeljenja riječi).
\newcolumntype{L}[1]{>{\raggedright\arraybackslash}p{#1}}
\newcommand{\codefile}[1]{\texttt{#1}}
\newcommand{\kbd}[1]{\texttt{#1}}

% ---------- Code frame environment ----------
% Automatski dodaje fragile i shrink=5, koristi se umjesto \begin{frame}
% Primjer: \begin{codeframe}{Naslov} ... \end{codeframe}
% Primjer s opcijama: \begin{codeframe}[shrink=15]{Naslov} ... \end{codeframe}
\newenvironment{codeframe}[2][]
  {\begin{frame}[fragile,shrink=5,#1]{#2}}
  {\end{frame}}

% ============================================================

\documentclass[aspectratio=149,11pt]{beamer}
\usetheme{CambridgeUS}

\usepackage{fontspec}              % radi samo sa XeLaTeX/LuaLaTeX
\usepackage[croatian]{babel}
\usepackage{csquotes}              % \enquote{} → hrvatski navodnici „..."
\usepackage{booktabs}
\usepackage{array}
\usepackage{listings}
\usepackage{xcolor}
\usepackage{colortbl}
\usepackage{hyperref}
\usepackage{tikz}
\usetikzlibrary{positioning,arrows.meta}

% ---------- Tipografija: IBM Plex ----------
\setsansfont{IBMPlexSans}[
  Extension      = .otf,
  UprightFont    = *-Regular,
  ItalicFont     = *-Italic,
  BoldFont       = *-SemiBold,
  BoldItalicFont = *-SemiBoldItalic
]
\setmainfont{IBMPlexSans}[
  Extension      = .otf,
  UprightFont    = *-Regular,
  ItalicFont     = *-Italic,
  BoldFont       = *-SemiBold,
  BoldItalicFont = *-SemiBoldItalic
]
\setmonofont{IBMPlexMono}[
  Extension      = .otf,
  UprightFont    = *-Regular,
  ItalicFont     = *-Italic,
  BoldFont       = *-SemiBold,
  BoldItalicFont = *-SemiBoldItalic,
  Scale          = 0.92
]

% ---------- Paleta ----------
\definecolor{ink}{HTML}{0C4685}      % primarna tamnoplava
\definecolor{ink2}{HTML}{6E7BAA}     % sekundarna plavo-siva
\definecolor{tint}{HTML}{EAEEF6}     % svjetli plavi tint
\definecolor{textclr}{HTML}{1A222E}  % glavni tekst
\definecolor{mutedclr}{HTML}{6B7484} % podnožje, sekundarni tekst
\definecolor{bgclr}{HTML}{F6F7F9}    % pozadina slajda
\definecolor{lineclr}{HTML}{D9DEE7}  % tanke linije tablica
\definecolor{divrule}{HTML}{9DB4D6}  % akcent na tamnim sekcijskim slajdovima
\definecolor{divtext}{HTML}{E6ECF6}  % tekst na tamnim sekcijskim slajdovima

\setbeamercolor{normal text}{fg=textclr,bg=bgclr}
\setbeamercolor{structure}{fg=ink}
\setbeamercolor{alerted text}{fg=ink}

\setbeamercolor{title}{fg=ink,bg=tint}
\setbeamercolor{subtitle}{fg=mutedclr}
\setbeamercolor{author}{fg=textclr}
\setbeamercolor{date}{fg=mutedclr}
\setbeamercolor{frametitle}{fg=ink,bg=tint}

\setbeamercolor{section in head/foot}{fg=white,bg=ink}
\setbeamercolor{subsection in head/foot}{fg=ink,bg=tint}
\setbeamercolor{author in head/foot}{fg=white,bg=ink}
\setbeamercolor{title in head/foot}{fg=textclr,bg=tint}
\setbeamercolor{date in head/foot}{fg=mutedclr,bg=tint}

\setbeamercolor{block title}{fg=white,bg=ink}
\setbeamercolor{block body}{fg=textclr,bg=tint}
\setbeamercolor{block title example}{fg=white,bg=ink2}
\setbeamercolor{block body example}{fg=textclr,bg=tint}
\setbeamercolor{itemize item}{fg=ink2}
\setbeamercolor{itemize subitem}{fg=ink2}
\setbeamercolor{enumerate item}{fg=ink2}

\setbeamertemplate{itemize items}[square]
\setbeamertemplate{navigation symbols}{}
\arrayrulecolor{lineclr}

% ---------- Kolegij, autori i logo ----------
\newcommand{\kolegij}{Programiranje za web}
\author{doc.\,dr.\,sc. Tomislav Petković \and mag.\,ing. Josip Torić}
\newcommand{\znak}[2]{\raisebox{-0.5\height}{\includegraphics[height=#1]{#2}}}
\titlegraphic{\znak{2.1cm}{../common/SIT_logo_tekst1.pdf}\hspace{1.8em}\znak{1.8cm}{../common/UNIZD_logo.pdf}}

% ---------- Zaglavlje ----------
% Jedna puna traka: lijevo sekcija, desno podsekcija.
\setbeamertemplate{headline}{%
  \leavevmode%
  \hbox{%
    \begin{beamercolorbox}[wd=\paperwidth,ht=2.25ex,dp=1ex,leftskip=.3cm,rightskip=.3cm]{section in head/foot}%
      \usebeamerfont{section in head/foot}\insertsectionhead\hfill\insertsubsectionhead%
    \end{beamercolorbox}%
  }%
}

% ---------- Podnožje ----------
% Lijevo: naziv trenutne prezentacije (kratki oblik naslova iz \title[...]{...})
% Sredina: naziv kolegija.  Desno: datum i redni broj slajda.
\setbeamertemplate{footline}{%
  \leavevmode%
  \hbox{%
    \begin{beamercolorbox}[wd=.34\paperwidth,ht=2.25ex,dp=1ex,center]{author in head/foot}%
      \usebeamerfont{author in head/foot}\insertshorttitle%
    \end{beamercolorbox}%
    \begin{beamercolorbox}[wd=.33\paperwidth,ht=2.25ex,dp=1ex,center]{title in head/foot}%
      \usebeamerfont{title in head/foot}\kolegij%
    \end{beamercolorbox}%
    \begin{beamercolorbox}[wd=.33\paperwidth,ht=2.25ex,dp=1ex,center]{date in head/foot}%
      \usebeamerfont{date in head/foot}\insertshortdate{}\hspace*{2em}%
      \insertframenumber{}\,/\,\inserttotalframenumber\hspace*{1ex}%
    \end{beamercolorbox}%
  }%
  \vskip0pt%
}

% ---------- Sekcijski slajd ----------
% \sekcija{Naslov sekcije}{jedna linija koja nabraja teme sekcije}
\newcommand{\sekcija}[2]{%
  \section{#1}%
  {%
    \setbeamercolor{background canvas}{bg=ink}%
    \setbeamertemplate{headline}{}%
    \begin{frame}
      \vfill
      \hspace*{0.03\textwidth}{\color{divrule}\rule{2.6cm}{1.1pt}}\par
      \vskip 0.7em
      \hspace*{0.03\textwidth}{\Large\bfseries\color{white}#1}\par
      \vskip 0.6em
      \hspace*{0.03\textwidth}\begin{minipage}{0.9\textwidth}\color{divtext}\small #2\end{minipage}\par
      \vfill
    \end{frame}%
  }%
}

% ---------- Listings ----------
\lstdefinelanguage{bash}{
  morekeywords={sudo,apt,brew,pip,pipx,uv,winget,python,git,flask,source,deactivate,code,gunicorn,docker,compose,heroku,render,railway,node,npm,pnpm,yarn,uvicorn,alembic,pytest,ruff,black,precommit},
  sensitive=false,
  morecomment=[l]{\#}
}
\lstdefinelanguage{python}{
  morekeywords={from,import,def,class,return,if,elif,else,for,in,while,try,except,raise,with,as,not,and,or,is,lambda,yield,True,False,None,await,async},
  sensitive=true,
  morecomment=[l]{\#},
  morestring=[b]',
  morestring=[b]"
}
\lstdefinelanguage{sql}{
  morekeywords={SELECT,FROM,WHERE,INSERT,INTO,VALUES,UPDATE,SET,DELETE,CREATE,TABLE,PRIMARY,KEY,FOREIGN,REFERENCES,UNIQUE,NOT,NULL,INDEX,ON,ALTER,DROP,CONSTRAINT},
  sensitive=false,
  morecomment=[l]{--}
}
\lstdefinelanguage{json}{
  morecomment=[l]{//},
  morestring=[b]",
  showstringspaces=false
}
\lstdefinelanguage{typescript}{
  morekeywords={import,export,from,const,let,var,function,return,if,else,for,of,in,while,switch,case,break,async,await,try,catch,finally,throw,new,class,this,extends,implements,interface,type,enum,readonly,public,private,protected,static,get,set,true,false,null,undefined,void,string,number,boolean,any,unknown,never,typeof,keyof,as},
  sensitive=true,
  morecomment=[l]{//},
  morecomment=[s]{/*}{*/},
  morestring=[b]',
  morestring=[b]",
  morestring=[b]`
}
\lstalias{javascript}{typescript}
\lstset{
  backgroundcolor=\color{tint},
  basicstyle=\ttfamily\footnotesize,
  keywordstyle=\color{ink}\bfseries,
  commentstyle=\color{mutedclr},
  frame=single,
  rulecolor=\color{lineclr},
  frameround=tttt,
  breaklines=true,
  tabsize=2,
  showstringspaces=false,
  upquote=true,
  columns=fullflexible,
  escapeinside={<@}{@>}
}
% ---------- Dijakritici u isječcima koda ----------
% listings čita bajt po bajt, pa višebajtni znak premjesti: "početni" ispadne
% "čpoetni". literate to ne popravlja. Rješenje je izlaz u LaTeX: sve između
% <@ i @> ne prolazi kroz listings, nego se slaže kao običan tekst, a unutar
% okvira je font ionako ttfamily, pa se ništa ne razlikuje.
%
%   async def get(...):
%       club = await repo.get_by_id(db, club_id)   <@\kom{dohvat kluba}@>
%       raise AppError("not_found", <@"Klub nije pronađen"@>, 404)
%
% Graničnici <@ i @> odabrani su jer se ne pojavljuju ni u jednom jeziku
% koji se na predavanjima citira; (* bi se sudaralo s def f(*args).
\newcommand{\kom}[1]{{\color{mutedclr}\# #1}}
% Isto za TypeScript, Vue i CSS, kojima komentar počinje s //:
%   const lifters = useLiftersStore()   <@\komts{dijeljeni popis}@>
\newcommand{\komts}[1]{{\color{mutedclr}// #1}}

% ---------- Činjenica naspram konvencije ----------
% Sve što nije propisano standardom ili jezikom, nego je dogovor struke,
% ide u ovaj okvir, da student zna koje je vrste tvrdnja pred njim.
% \begin{konvencija} ... \end{konvencija}
% \begin{konvencija}[Konvencija: nazivi putanja] ... \end{konvencija}
\newenvironment{konvencija}[1][Konvencija]
  {\begin{exampleblock}{#1}}
  {\end{exampleblock}}

% ---------- Helper macros ----------
% Stupac tablice zadane širine, poravnat ulijevo (bez rastezanja i dijeljenja riječi).
\newcolumntype{L}[1]{>{\raggedright\arraybackslash}p{#1}}
\newcommand{\codefile}[1]{\texttt{#1}}
\newcommand{\kbd}[1]{\texttt{#1}}

% ---------- Code frame environment ----------
% Automatski dodaje fragile i shrink=5, koristi se umjesto \begin{frame}
% Primjer: \begin{codeframe}{Naslov} ... \end{codeframe}
% Primjer s opcijama: \begin{codeframe}[shrink=15]{Naslov} ... \end{codeframe}
\newenvironment{codeframe}[2][]
  {\begin{frame}[fragile,shrink=5,#1]{#2}}
  {\end{frame}}


\title[Predavanje 5: Oblikovanje API-ja]{Programiranje za web \\ Predavanje 5: Oblikovanje API-ja i CRUD domene}
\date{05.11.2026.}

\begin{document}

\frame{\titlepage}

% ------------------------------------------------------------
\begin{frame}{Plan predavanja}
\begin{enumerate}
  \item Od zapisa do resursa: adrese, operacije i prava
  \item Svojstva HTTP metoda: idempotentnost i izbor metode
  \item Sheme: što ulazi, što izlazi i što se provjerava
  \item Popisi: granica odgovora, straničenje i filtriranje
\end{enumerate}

\vspace{0.3em}
Dosad je aplikacija znala tko je korisnik i što smije. Sada dobiva zapise nad kojima radi i API kojim ih netko drugi može stvarati, čitati, mijenjati i brisati.
\end{frame}

% ------------------------------------------------------------
\begin{frame}{Ishodi učenja}
Nakon ovog predavanja student može:

\begin{itemize}
  \item preslikati četiri osnovne operacije nad zapisom u metode i statusne kodove
  \item obrazložiti izbor između ugniježđene i samostalne rute
  \item napisati odvojene sheme za unos, izmjenu i odgovor i reći zašto se razlikuju
  \item dodati provjeru koja nadilazi tip polja i reći gdje je granica sheme
  \item ograničiti veličinu odgovora straničenjem i filtriranjem kroz parametre upita
  \item objasniti zašto se povezani zapisi moraju zatražiti unaprijed
\end{itemize}
\end{frame}

% ------------------------------------------------------------
\begin{frame}{Domena primjera}
Sustav prijava klubova na powerlifting natjecanje. Uz klub i račun dolaze tri nova zapisa:

\vspace{0.2em}
\begin{center}
\begin{tikzpicture}[
  font=\small,
  ent/.style={draw=ink, thick, rounded corners, fill=white, inner sep=3mm, text=textclr, minimum width=26mm},
  lbl/.style={font=\scriptsize, text=mutedclr}
]
  \node[ent] (club) {\textbf{Klub}};
  \node[ent, right=24mm of club] (lifter) {\textbf{Natjecatelj}};
  \node[ent, below=14mm of lifter] (reg) {\textbf{Prijava}};
  \node[ent, left=24mm of reg] (comp) {\textbf{Natjecanje}};

  \draw[-{Latex}, ink2, thick] (lifter.west) -- node[lbl, above] {1:N} (club.east);
  \draw[-{Latex}, ink2, thick] (reg.north) -- node[lbl, right] {1:N} (lifter.south);
  \draw[-{Latex}, ink2, thick] (reg.west) -- node[lbl, above] {1:N} (comp.east);
\end{tikzpicture}
\end{center}

\vspace{0.4em}
Prijava nije samo spojnica dvaju zapisa. Nosi i vlastite podatke, težinsku kategoriju i najbolji zbroj (\texttt{total}), pa je i sama zapis o kojem se vodi računa.
\end{frame}

% ============================================================
\sekcija{Od zapisa do resursa}{Što je resurs, četiri operacije nad njim, pripadnost u adresi, prava koja ruta objavljuje.}

% ------------------------------------------------------------
\begin{frame}{Zapis i resurs}
Zapis u bazi nema adresu. Postoji kao redak, a do njega se dolazi upitom koji zna kako je tablica složena. Klijent tu bazu nema: raspolaže samo adresom i metodom.

\vspace{0.9em}
\begin{center}
\begin{tikzpicture}[
  kut/.style={draw=ink, thick, rounded corners, fill=white, align=left,
              text width=44mm, font=\footnotesize, text=textclr, inner sep=3mm},
  vanj/.style={draw=lineclr, thick, dashed, rounded corners, fill=bgclr, align=left,
              text width=44mm, font=\footnotesize, text=mutedclr, inner sep=3mm},
  str/.style={-{Latex}, ink2, thick},
  opis/.style={font=\scriptsize, text=mutedclr}
]
  \node[kut] (b) {\textbf{tablica \texttt{lifters}}\\[3pt]
    \texttt{31\quad Ana\quad Horvat}\\
    \texttt{32\quad Ivan\quad Perić}};
  \node[vanj, right=22mm of b] (k) {\textbf{što vidi klijent}\\[3pt]
    \texttt{GET /lifters}\\
    \texttt{GET /lifters/31}};
  \draw[str] (b) -- (k);
  \node[opis, below=2mm of b] {redak, stupac, upit};
  \node[opis, below=2mm of k] {adresa i metoda};
\end{tikzpicture}
\end{center}

\vspace{0.9em}
Zato se svaka vrsta zapisa objavljuje kao \textbf{resurs}: nešto što ima svoju adresu i skup radnji koje se nad njom izvode.
\end{frame}

% ------------------------------------------------------------
\begin{frame}{Odnos resursa i tablice}
\vspace{0.4em}
\begin{center}
{\small
\setlength{\tabcolsep}{10pt}
\renewcommand{\arraystretch}{1.3}
\begin{tabular}{@{}ll@{}}
  \toprule
  \textbf{U bazi} & \textbf{Prema klijentu} \\
  \midrule
  tablica \texttt{lifters}    & resurs na \texttt{/lifters} \\
  jedan redak                 & \texttt{/lifters/31} \\
  stupac \texttt{first\_name} & polje u tijelu zahtjeva i odgovora \\
  \bottomrule
\end{tabular}
}
\end{center}

\vspace{1.1em}
Preslikavanje ne mora biti jedan na jedan. Resurs smije spojiti podatke iz više tablica, izostaviti stupce koje nitko izvana ne treba, ili uopće ne odgovarati nijednoj tablici.
\end{frame}

% ------------------------------------------------------------
\begin{frame}{Podsjetnik: statusni kodovi}
Operacije se razlikuju i po kodu kojim završavaju, pa se isplati obnoviti ono iz prvoga predavanja. Prva znamenka nosi razred:

\vspace{0.6em}
\begin{center}
{\small
\setlength{\tabcolsep}{10pt}
\renewcommand{\arraystretch}{1.25}
\begin{tabular}{@{}ll@{}}
  \toprule
  \textbf{Kod} & \textbf{Značenje} \\
  \midrule
  \texttt{200} & obrađeno, tijelo slijedi \\
  \texttt{201} & nastao je novi zapis \\
  \texttt{204} & obrađeno, tijela nema \\
  \texttt{404} & na toj putanji nema ničega \\
  \texttt{409} & sukob sa stanjem resursa \\
  \texttt{422} & sadržaj se ne može obraditi \\
  \bottomrule
\end{tabular}
}
\end{center}

\vspace{0.6em}
Svi iz razreda \texttt{2xx} znače da je zahtjev obrađen, a oni iz \texttt{4xx} da s njim nešto nije u redu.
\end{frame}

% ------------------------------------------------------------
\begin{frame}{Četiri operacije nad zapisom}
Nad svakim se resursom radi isto: stvara se, čita, mijenja i briše. Te operacije nose naziv CRUD (engl. \emph{create, read, update, delete}):

\vspace{1.1em}
\begin{center}
{\small
\setlength{\tabcolsep}{10pt}
\renewcommand{\arraystretch}{1.3}
\begin{tabular}{@{}llll@{}}
  \toprule
  \textbf{Operacija} & \textbf{Metoda} & \textbf{Kod} & \textbf{Tijelo odgovora} \\
  \midrule
  popis        & \texttt{GET}    & \texttt{200} & niz zapisa \\
  jedan zapis  & \texttt{GET}    & \texttt{200} & zapis \\
  stvaranje    & \texttt{POST}   & \texttt{201} & stvoreni zapis \\
  izmjena      & \texttt{PATCH}  & \texttt{200} & izmijenjeni zapis \\
  brisanje     & \texttt{DELETE} & \texttt{204} & / \\
  \bottomrule
\end{tabular}
}
\end{center}
\end{frame}

% ------------------------------------------------------------
\begin{frame}{Tijelo odgovora uz \texttt{201} i \texttt{204}}
Dva retka iz tablice odstupaju:

\vspace{0.9em}
\begin{itemize}
  \itemsep0.6em
  \item \textbf{Stvaranje vraća \texttt{201} i zapis.} Klijent prije odgovora ne zna oznaku koju dodjeljuje baza, pa je mora dobiti natrag; bez nje ne može uputiti nijedan sljedeći zahtjev nad tim zapisom.
  \item \textbf{Brisanje vraća \texttt{204} i ništa.} Zapisa više nema, pa nema ni što poslati. Tijelo uz \texttt{204} standard i ne dopušta.
\end{itemize}

\begin{konvencija}
Sama preslika operacija u metode nije propisana. \texttt{POST} i \texttt{GET} standard opisuje, ali nigdje ne piše da izmjena mora biti \texttt{PATCH} ni da brisanje mora vratiti \texttt{204}. To je dogovor struke. Korist od njega je u tome što ga se držimo na svim resursima, pa tko nauči jedan, zna i ostale.
\end{konvencija}
\end{frame}

% ------------------------------------------------------------
\begin{frame}{Pripadnost u adresi}
Natjecatelj uvijek pripada nekom klubu. Ta se pripadnost u adresi može izraziti na dva načina:

\vspace{0.25em}
\begin{center}
{\small
\setlength{\tabcolsep}{8pt}
\renewcommand{\arraystretch}{1.05}
\begin{tabular}{@{}ll@{}}
  \toprule
  \textbf{Oblik} & \textbf{Adresa} \\
  \midrule
  ugniježđena & \texttt{/clubs/12/lifters} \\
  samostalna  & \texttt{/lifters?club\_id=12} \\
  \bottomrule
\end{tabular}
}
\end{center}

\vspace{0.3em}
U prvom je obliku pripadnost dio adrese i ne može se izostaviti. U drugom je uvjet pretrage: izostavi li se, ruta vraća natjecatelje svih klubova.
\begin{konvencija}[Izbor, ne pravilo]
Ugniježđeni oblik odgovara odnosu u kojem podređeni zapis bez nadređenoga nema smisla. Samostalni je prikladniji kad se zapis dohvaća po više kriterija. U našem su projektu natjecatelji ugniježđeni pod klubom, a natjecanja stoje samostalno.
\end{konvencija}
\end{frame}

% ------------------------------------------------------------
\begin{frame}[fragile]{Mjesto ugnježđivanja}
Putanje u datoteci s rutama natjecatelja zapisane su bez prefiksa: \texttt{/} za popis i stvaranje, \texttt{/\{lifter\_id\}} za pojedinačni zapis. Puna adresa nastaje tek pri uključivanju u aplikaciju, nadovezivanjem prefiksa na te putanje:

\begin{lstlisting}[language=python]
# app/main.py
app.include_router(clubs_router, prefix="/clubs")
app.include_router(lifters_router,
                   prefix="/clubs/{club_id}/lifters")
app.include_router(competitions_router, prefix="/competitions")
app.include_router(registrations_router,
                   prefix="/competitions/{comp_id}/registrations")
\end{lstlisting}

\vspace{0.4em}
Time je oblik API-ja vidljiv na jednom mjestu. Iz ovog se popisa čita koja je pripadnost ugniježđena, bez otvaranja ijedne datoteke s rutama.

\vspace{0.9em}
Oznaka \texttt{\{club\_id\}} time stoji u prefiksu, a ne u putanji same rute.
\end{frame}

% ------------------------------------------------------------
\begin{frame}[fragile]{Oznaka iz prefiksa}
\begin{lstlisting}[language=python]
# app/routers/lifters.py
@router.get("/", response_model=list[LifterResponse])
async def list_lifters(
    club_id: int,
    db: AsyncSession = Depends(get_db),
    user: User = Depends(require_role("admin", "club")),
):
    return await lifter_service.list_lifters(db, club_id, user, ...)
\end{lstlisting}

\vspace{0.4em}
Framework traži \texttt{club\_id} u cijeloj adresi, dakle i u prefiksu, i preda ga kao običan argument.

\vspace{0.9em}
Ruta time ne zna ni da je ugniježđena. Isti kod radi i da je adresa \texttt{/lifters/\{club\_id\}}.
\end{frame}

% ------------------------------------------------------------
\begin{frame}[fragile]{Različita prava za čitanje i pisanje}
Natjecanja čita svatko prijavljen, a mijenja ih samo organizator:

\begin{lstlisting}[language=python]
# app/routers/competitions.py
@router.get("/", response_model=list[CompetitionResponse])
async def list_competitions(
    db: AsyncSession = Depends(get_db),
    user: User = Depends(get_current_user),
):
    return await competition_service.list_competitions(db)

@router.post("/", response_model=CompetitionResponse, status_code=201)
async def create_competition(
    body: CompetitionCreate,
    db: AsyncSession = Depends(get_db),
    admin: User = Depends(require_role("admin")),
):
    return await competition_service.create_competition(db, body)
\end{lstlisting}
\end{frame}

% ------------------------------------------------------------
\begin{frame}{Uloga u potpisu rute}
\texttt{get\_current\_user} traži samo to da je korisnik utvrđen, bez obzira na to koja mu je uloga. \texttt{require\_role("admin")} uz to traži i određenu ulogu.

\vspace{1.1em}
Tko rutu smije pozvati time se čita iz njezina potpisa, a ne iz tijela funkcije. Ni \texttt{list\_competitions} ni \texttt{create\_competition} o pravima ne odlučuju same.

\vspace{1.1em}
Uz \texttt{POST} stoji i \texttt{status\_code=201}, jer stvaranje vraća drukčiji kod od ostalih operacija.
\end{frame}

% ------------------------------------------------------------
\begin{frame}{Brisanje i povezani zapisi}
Obriše li se natjecatelj, prijave koje na njega pokazuju ostaju bez zapisa na koji se odnose. Tri su izlaza iz toga:

\vspace{1.1em}
\begin{itemize}
  \itemsep0.7em
  \item \textbf{Zabraniti dok veze postoje.} Zahtjev se odbija, zapis ostaje netaknut.
  \item \textbf{Obrisati i povezane zapise.} Nestaje i prijava i svaki trag o njoj.
  \item \textbf{Označiti stanje umjesto brisanja.} Zapis ostaje, a stupac kaže da više ne vrijedi.
\end{itemize}

\vspace{1.1em}
Izbor ovisi o tome treba li povijest ostati dostupna.
\end{frame}

% ------------------------------------------------------------
\begin{frame}{Izostanak pravila o brisanju}
Da izbora nema. Ruta koja briše bez ijednog od tih pravila završava neuhvaćenom pogreškom baze umjesto odgovorom koji klijent razumije.

\begin{konvencija}[Izbor, ne pravilo]
U ovom projektu natjecatelj se ne briše dok ima prijave: zahtjev vraća \texttt{409} i kod \texttt{has\_registrations}. Prijava se ne briše nego joj se \texttt{status} postavlja na \texttt{"withdrawn"}, pa odjava ostaje vidljiva.

\vspace{0.5em}
Brisanje povezanih zapisa nigdje ne upotrebljavamo, jer povijest prijava na natjecanje čuvamo.
\end{konvencija}
\end{frame}

% ------------------------------------------------------------
\begin{frame}{Vježba: četiri operacije u Swaggeru}
Prijavite se u \texttt{/docs} kao klub \texttt{behemot@plreg.hr} i radite nad njegovim natjecateljima.

\vspace{1.1em}
\begin{enumerate}
  \itemsep0.6em
  \item Stvorite natjecatelja, pa nad istim zapisom izvedite \texttt{GET}, \texttt{PATCH} i \texttt{DELETE}. Za svaki zabilježite kod odgovora i ima li tijela.
  \item Stvorite još jednoga, prijavite ga na \emph{Zimski kup}, pa ga pokušajte obrisati.
\end{enumerate}

\vspace{1.1em}
Koji ste kod dobili u drugom koraku, tko ga je podigao i zašto?

\vspace{0.7em}
Po čemu se sve razlikuju te četiri operacije? (adresa, tijelo zahtjeva, učinak na stanje, ponovljivost)
\end{frame}

% ============================================================
\sekcija{Svojstva HTTP metoda}{Podsjetnik na metode, idempotentnost, razlika između \texttt{PUT}-a i \texttt{PATCH}-a, izbor metode za izmjenu.}

% ------------------------------------------------------------
\begin{frame}{Podsjetnik: HTTP metode}
\vspace{0.3em}
\begin{center}
{\small
\setlength{\tabcolsep}{10pt}
\renewcommand{\arraystretch}{1.1}
\begin{tabular}{@{}lL{7.6cm}@{}}
  \toprule
  \textbf{Metoda} & \textbf{Što traži od poslužitelja} \\
  \midrule
  \texttt{GET}     & prikaz resursa \\
  \texttt{HEAD}    & isto što i \texttt{GET}, ali bez tijela odgovora \\
  \texttt{OPTIONS} & popis metoda koje resurs podržava \\
  \texttt{TRACE}   & zahtjev vraćen onakav kakav je stigao, za dijagnostiku \\
  \texttt{POST}    & obradu poslanog sadržaja, najčešće stvaranje zapisa \\
  \texttt{PUT}     & da resurs postane poslani sadržaj \\
  \texttt{PATCH}   & izmjenu dijela resursa \\
  \texttt{DELETE}  & brisanje resursa \\
  \bottomrule
\end{tabular}
}
\end{center}

\end{frame}

% ------------------------------------------------------------
\begin{frame}{Idempotentnost}
Metoda je idempotentna ako niz od više jednakih zahtjeva ostavlja poslužitelj u istom stanju kao i jedan takav zahtjev. Standard to svojstvo pridružuje samim metodama. Za koje od njih vrijedi?

\vspace{0.6em}
\begin{center}
{\small
\setlength{\tabcolsep}{14pt}
\renewcommand{\arraystretch}{1.15}
\begin{tabular}{@{}ll@{}}
  \toprule
  \textbf{Metoda} & \textbf{Idempotentna} \\
  \midrule
  \texttt{GET}     & \\
  \texttt{HEAD}    & \\
  \texttt{OPTIONS} & \\
  \texttt{TRACE}   & \\
  \texttt{POST}    & \\
  \texttt{PUT}     & \\
  \texttt{PATCH}   & \\
  \texttt{DELETE}  & \\
  \bottomrule
\end{tabular}
}
\end{center}
\end{frame}

% ------------------------------------------------------------
\begin{frame}{Idempotentne metode}
Metoda je idempotentna ako niz od više jednakih zahtjeva ostavlja poslužitelj u istom stanju kao i jedan takav zahtjev.

\vspace{0.6em}
\begin{center}
{\small
\setlength{\tabcolsep}{14pt}
\renewcommand{\arraystretch}{1.15}
\begin{tabular}{@{}ll@{}}
  \toprule
  \textbf{Metoda} & \textbf{Idempotentna} \\
  \midrule
  \texttt{GET}     & da \\
  \texttt{HEAD}    & da \\
  \texttt{OPTIONS} & da \\
  \texttt{TRACE}   & da \\
  \texttt{POST}    & \textbf{ne} \\
  \texttt{PUT}     & da \\
  \texttt{PATCH}   & \textbf{ne} \\
  \texttt{DELETE}  & da \\
  \bottomrule
\end{tabular}
}
\end{center}
\end{frame}

% ------------------------------------------------------------
\begin{frame}{Posljedice idempotentnosti}
\begin{itemize}
  \itemsep0.6em
  \item \texttt{GET}, \texttt{HEAD}, \texttt{OPTIONS} i \texttt{TRACE} ne mijenjaju stanje.
  \item \texttt{PUT} i \texttt{DELETE} mijenjaju stanje, ali ponavljanje na njemu ništa ne mijenja: zapis se postavlja na poslani sadržaj, odnosno briše.
  \item \texttt{POST} i \texttt{PATCH} nisu idempotentni.
\end{itemize}

\vspace{1.1em}
Klijent zato smije ponoviti neuspjeli \texttt{GET} ili \texttt{DELETE} bez razmišljanja.

\vspace{0.7em}
Ponovljeni \texttt{POST} može stvoriti drugi zapis, pa zaštita od toga mora doći iz aplikacije, najčešće ograničenjem jedinstvenosti u bazi.
\end{frame}

% ------------------------------------------------------------

\begin{frame}{\texttt{PUT} i \texttt{PATCH}}
Obje mijenjaju postojeći zapis, ali tijelo zahtjeva kod njih ne znači isto. \texttt{PUT} šalje ciljno stanje zapisa, \texttt{PATCH} opis izmjene nad onim što već postoji.

\vspace{0.9em}
\begin{center}
{\small
\setlength{\tabcolsep}{10pt}
\renewcommand{\arraystretch}{1.3}
\begin{tabular}{@{}lll@{}}
  \toprule
   & \textbf{\texttt{PUT}} & \textbf{\texttt{PATCH}} \\
  \midrule
  Tijelo zahtjeva     & ciljno stanje & opis izmjene \\
  Izostavljeno polje  & briše se      & ne mijenja se \\
  Ponovljeni zahtjev  & isto stanje   & ne nužno isto stanje \\
  Obrada u aplikaciji & zamjena       & spajanje s postojećim \\
  \bottomrule
\end{tabular}
}
\end{center}

\vspace{0.9em}
Ostala tri retka slijede iz prvoga. Polje koje klijent zaboravi poslati \texttt{PUT} briše, jer ga u ciljnom stanju nema. Ponavljanje \texttt{PATCH}-a ovisi o tome kakvu izmjenu nosi: \texttt{\{"first\_name": "Ana"\}} ga podnosi, \texttt{\{"total": "+5"\}} ne.
\end{frame}

% ------------------------------------------------------------
\begin{frame}{Odabir metode za izmjenu}
\textbf{\texttt{PUT}} odgovara kad klijent drži cjelovit prikaz zapisa i mijenja ga u cijelosti, primjerice postavke ili dokument. Tada je brisanje polja izričit postupak: polje koje klijent više ne pošalje nestaje. Obrada je jednostavnija, jer nema spajanja ni dvoznačnosti između neposlanog polja i poslanog \texttt{null}.

\vspace{0.25em}
\textbf{\texttt{PATCH}} odgovara kad je zapis velik, a izmjena malena, i kad klijent uopće ne poznaje sva polja. Zamjena bi tada izbrisala sve čega se klijent ne sjeća poslati, pa je slanje samo izmijenjenog i sigurnije i jeftinije.

\vspace{0.25em}
\begin{konvencija}[Izbor, ne pravilo]
U našem projektu izmjena je svugdje izvedena metodom \texttt{PATCH}, jer klub mijenja jedno ili dva polja natjecatelja i ne drži cijeli zapis. API koji se odluči za \texttt{PUT} ništa ne krši, ali mora računati na to da izostavljeno polje znači brisanje.
\end{konvencija}
\end{frame}

% ============================================================
\sekcija{Sheme}{Što shema provjerava, zašto jedna nije dovoljna, unos i odgovor, provjere koje nadilaze tip, izmjena i neposlana polja.}

% ------------------------------------------------------------
\begin{frame}{Shema kao opis ulaza i izlaza}
Shema je zapis o tome kakve podatke ruta prima i kakve vraća. Framework je provjerava prije nego zahtjev dođe do poslovnih pravila, pa neispravan sadržaj do njih ne stiže.

\vspace{1.1em}
Shema provjerava troje:

\vspace{0.8em}
\begin{itemize}
  \itemsep0.7em
  \item jesu li obavezna polja prisutna
  \item je li svako polje očekivanog tipa
  \item zadovoljava li vrijednost ograničenje zapisano uz polje, primjerice popis dopuštenih vrijednosti ili raspon
\end{itemize}
\end{frame}

% ------------------------------------------------------------
\begin{frame}[fragile]{Shema: primjer}
\begin{lstlisting}[language=python]
# app/schemas/lifter.py
class LifterCreate(BaseModel):
    first_name: str = Field(min_length=1, max_length=80)
    birth_date: date
    gender: Literal["M", "F"]
\end{lstlisting}

\vspace{0.4em}
{\small
\begin{tabular}{@{}lL{7.2cm}@{}}
  \toprule
  \textbf{Poslano} & \textbf{Ishod} \\
  \midrule
  \texttt{"gender": "M"}   & prolazi dalje \\
  \texttt{"gender": "X"}   & odbijeno, vrijednost nije na popisu \\
  \texttt{"gender": 3}     & odbijeno; broj ne bi prošao ni da popisa nema \\
  \texttt{"first\_name": ""} & odbijeno, kraće od najmanje duljine \\
  bez \texttt{birth\_date}  & odbijeno, obavezno polje nedostaje \\
  \bottomrule
\end{tabular}
}

\vspace{0.4em}
\texttt{Literal} nabraja dopuštene vrijednosti, a \texttt{Field} ograničava duljinu: obje provjere stoje uz samo polje.
\end{frame}

% ------------------------------------------------------------

\begin{frame}{Isti zapis, različit skup polja}
Natjecatelj je u bazi jedan zapis, ali za svaku vrstu zahtjeva putuje drukčiji skup polja. Polje koje se već nalazi u adresi rute u tijelu se ne ponavlja:

\vspace{0.3em}
\begin{center}
{\small
\setlength{\tabcolsep}{0pt}
\renewcommand{\arraystretch}{1.18}
\begin{tabular}{@{}L{3.3cm}L{2.7cm}L{3.0cm}L{2.0cm}@{}}
  \toprule
  \textbf{Polje} & \textbf{Stvaranje} & \textbf{Izmjena} & \textbf{Odgovor} \\
  \midrule
  \texttt{id}       & nema ga   & u adresi       & uvijek \\
  ime i prezime     & obavezno  & ako se mijenja & uvijek \\
  datum rođenja     & obavezno  & ako se mijenja & uvijek \\
  spol              & obavezno  & ako se mijenja & uvijek \\
  \texttt{club\_id} & u adresi  & u adresi       & uvijek \\
  \bottomrule
\end{tabular}
}
\end{center}

\vspace{0.3em}
Stvaranje i izmjena traže različito, a odgovor nosi sve.
\end{frame}

% ------------------------------------------------------------
\begin{frame}{Odakle dolazi koje polje}
Šest polja u odgovoru stiže trima različitim putovima.

\vspace{0.3em}
\begin{center}
\begin{tikzpicture}[
  node distance=4mm,
  kut/.style={draw=ink, thick, rounded corners, fill=white, align=left,
              text width=46mm, font=\footnotesize, text=textclr, inner sep=2mm},
  vanj/.style={draw=lineclr, thick, dashed, rounded corners, fill=bgclr, align=left,
               text width=46mm, font=\footnotesize, text=mutedclr, inner sep=2mm},
  odg/.style={draw=ink, thick, rounded corners, fill=white, align=left,
              text width=38mm, font=\footnotesize, text=textclr, inner sep=2.5mm},
  str/.style={-{Latex}, ink2, thick}
]
  \node[kut] (a) {\textbf{adresa rute}\\klijent šalje \texttt{club\_id}};
  \node[kut, below=of a] (t) {\textbf{tijelo zahtjeva}\\klijent šalje ime, prezime,\\datum i spol};
  \node[vanj, below=of t] (b) {\textbf{baza}\\dodjeljuje \texttt{id}};
  \node[odg, right=14mm of t] (o) {\textbf{odgovor}\\\texttt{id}, ime, prezime,\\datum, spol, \texttt{club\_id}};

  \draw[str] (a.east) -- (o.west);
  \draw[str] (t.east) -- (o.west);
  \draw[str] (b.east) -- (o.west);
\end{tikzpicture}
\end{center}

\vspace{0.3em}
Klijent šalje pet od šest polja, ali samo četiri u tijelu. \texttt{club\_id} putuje adresom, a \texttt{id} ne dolazi od njega. Zato ih \texttt{LifterCreate} nema ni jedno.
\end{frame}

% ------------------------------------------------------------
\begin{frame}{Tri uloge istoga polja}
Jedna shema ne bi mogla ispuniti sve tri uloge, jer bi ista polja morala istodobno tražiti i ne tražiti.

\vspace{1.1em}
\begin{itemize}
  \itemsep0.7em
  \item Pri \textbf{stvaranju} datum rođenja mora biti prisutan.
  \item Pri \textbf{izmjeni} isti datum smije izostati, jer se mijenja samo ime.
  \item U \textbf{odgovoru} uvijek postoji, zajedno s \texttt{id} i \texttt{club\_id}, koje klijent nikada ne šalje.
\end{itemize}

\vspace{1.1em}
Zato ih razdvajamo: \texttt{LifterCreate}, \texttt{LifterUpdate} i \texttt{LifterResponse}.
\end{frame}

% ------------------------------------------------------------
\begin{frame}{Broj shema po zapisu}
Broj shema slijedi iz broja oblika u kojima se zapis pojavljuje:

\vspace{1.1em}
\begin{itemize}
  \itemsep0.7em
  \item \textbf{Autentikacija ima posve drukčiji skup.} Ondje stoje \texttt{LoginRequest}, \texttt{TokenResponse} i \texttt{RefreshRequest}, jer se ne stvara zapis nego izdaje token.
  \item \textbf{Zapis koji se nikada ne mijenja nema shemu za izmjenu.} Nedostaje joj uloga koju bi ispunjavala.
\end{itemize}

\begin{konvencija}
Podjela na unos, izmjenu i odgovor nije propisana nego uobičajena. Broj shema slijedi iz toga u koliko se oblika isti zapis pojavljuje.
\end{konvencija}
\end{frame}

% ------------------------------------------------------------
\begin{frame}[fragile]{Shema za unos}
Ruta koja stvara natjecatelja navodi shemu kao tip tijela:

\begin{lstlisting}[language=python]
# app/routers/lifters.py
async def create_lifter(club_id: int, body: LifterCreate, ...):
\end{lstlisting}

\vspace{0.6em}
\begin{lstlisting}[language=python]
# app/schemas/lifter.py
class LifterCreate(BaseModel):
    first_name: str = Field(min_length=1, max_length=80)
    last_name: str = Field(min_length=1, max_length=80)
    birth_date: date
    gender: Literal["M", "F"]
\end{lstlisting}

\vspace{0.4em}
Sve što u \texttt{LifterCreate} ne stoji, u tijelo zahtjeva ne ulazi.
\end{frame}

% ------------------------------------------------------------
\begin{frame}{Isti natjecatelj, dva oblika}
Redak u bazi i tijelo odgovora nisu ista stvar:

\vspace{0.7em}
\begin{center}
\begin{tikzpicture}[
  kut/.style={draw=ink, thick, rounded corners, fill=white, align=left,
              text width=44mm, font=\scriptsize\ttfamily, text=textclr, inner sep=2.5mm},
  str/.style={-{Latex}, ink2, thick},
  opis/.style={font=\scriptsize, text=mutedclr}
]
  \node[kut] (b) {id\quad\ \ 31\\club\_id\ 12\\first\_name\ \ Ana\\last\_name\ \ \ Horvat\\birth\_date\ \ 2001-05-02\\gender\quad\ \ F};
  \node[kut, right=32mm of b] (k) {\{\\\ \ "id": 31,\\\ \ "club\_id": 12,\\\ \ "first\_name": "Ana",\\\ \ ...\\\}};
  \draw[str] (b) -- node[opis, above, yshift=0.6mm] {\texttt{LifterResponse}} (k);
  \node[opis, below=2mm of b] {prema bazi: redak tablice};
  \node[opis, below=2mm of k] {prema klijentu: JSON};
\end{tikzpicture}
\end{center}

\vspace{0.7em}
Shema odgovora je pravilo pretvorbe između ta dva oblika.
\end{frame}

% ------------------------------------------------------------
\begin{frame}[fragile]{Zapis i shema odgovora}
\begin{lstlisting}[language=python]
# app/models/lifter.py
class Lifter(Base):
    __tablename__ = "lifters"
    id: Mapped[int] = mapped_column(primary_key=True)
    club_id: Mapped[int] = mapped_column(ForeignKey("clubs.id"))
    first_name: Mapped[str] = mapped_column(String(80))
    ...
\end{lstlisting}

\vspace{0.15em}
\begin{lstlisting}[language=python]
# app/schemas/lifter.py
class LifterResponse(BaseModel):
    id: int
    club_id: int
    first_name: str
    ...
\end{lstlisting}

\vspace{0.15em}
\codefile{repositories} vraća prvi oblik, \codefile{services} ga prosljeđuje, a ruta mora vratiti drugi.
\end{frame}

% ------------------------------------------------------------
\begin{frame}[fragile]{Prepisivanje polja rukom}
Pydantic po zadanom čita rječnike, ne atribute objekta, pa bi se polja prepisivala jedno po jedno:

\begin{lstlisting}[language=python]
# app/routers/lifters.py
return LifterResponse(
    id=lifter.id,
    first_name=lifter.first_name,
    last_name=lifter.last_name,
    birth_date=lifter.birth_date,
    gender=lifter.gender,
    club_id=lifter.club_id,
)
\end{lstlisting}

\vspace{0.4em}
Taj blok stajao bi uz svaku rutu koja vraća natjecatelja: popis, dohvat, stvaranje i izmjenu. Svako novo polje tražilo bi izmjenu na sva četiri mjesta.
\end{frame}

% ------------------------------------------------------------
\begin{frame}[fragile]{Jedan redak umjesto prepisivanja}
\begin{lstlisting}[language=python]
# app/schemas/lifter.py
class LifterResponse(BaseModel):
    id: int
    club_id: int
    first_name: str
    ...

    model_config = {"from_attributes": True}
\end{lstlisting}

\vspace{0.4em}
Zbog zadnjeg retka ruta smije vratiti \texttt{Lifter} izravno. Prepisivanje polja s prethodnog slajda otpada.

\vspace{0.7em}
Izlazna shema nije samo opis. Odgovor prolazi kroz nju, pa polje koje u njoj ne stoji ne izlazi ni greškom.
\end{frame}

% ------------------------------------------------------------
\begin{frame}[fragile]{Gdje se sheme navode}
Obje sheme stoje u potpisu rute, svaka na svojem mjestu:

\begin{lstlisting}[language=python]
# app/routers/lifters.py
@router.post("/", response_model=LifterResponse, status_code=201)
async def create_lifter(
    club_id: int,
    body: LifterCreate,
    db: AsyncSession = Depends(get_db),
    user: User = Depends(require_role("admin", "club")),
):
    return await lifter_service.create_lifter(db, club_id, body, user)
\end{lstlisting}

\vspace{0.4em}
\texttt{body: LifterCreate} opisuje što ulazi, pa framework tijelo zahtjeva provjeri prije nego funkcija počne raditi. \texttt{response\_model=LifterResponse} opisuje što izlazi, pa pretvorbu iz \texttt{Lifter} u odgovor obavlja framework.

\vspace{0.7em}
Iz ta dva podatka nastaje i opis rute u \texttt{/docs}.
\end{frame}

% ------------------------------------------------------------
\begin{frame}[fragile]{Provjere koje nadilaze tip}
Tip polja hvata pogrešan oblik, ali ne i besmislenu vrijednost. Datum rođenja u budućnosti je ispravan datum:

\begin{lstlisting}[language=python]
# app/schemas/lifter.py
class LifterCreate(BaseModel):
    birth_date: date

    @field_validator("birth_date")
    @classmethod
    def must_be_past(cls, v: date) -> date:
        if v >= date.today():
            raise ValueError(<@"Datum rođenja mora biti u prošlosti"@>)
        return v
\end{lstlisting}

\vspace{0.4em}
Kad takva provjera padne, odgovor ima isti oblik kao i pogreška u tipu: imenuje polje i navodi razlog. Sloj s poslovnim pravilima time nikada ne prima vrijednost koja nema smisla.
\end{frame}

% ------------------------------------------------------------
\begin{frame}[fragile]{Doseg sheme}
Shema vidi poslani sadržaj i današnji datum, pa može usporediti i dva polja istog zahtjeva:

\begin{lstlisting}[language=python]
# app/schemas/competition.py
@field_validator("final_deadline")
@classmethod
def final_after_prelim(cls, v, info):
    prelim = info.data.get("prelim_deadline")
    if prelim and v <= prelim:
        raise ValueError(<@"final\_deadline mora biti nakon prelim\_deadline"@>)
    return v
\end{lstlisting}

\vspace{0.4em}
Rokovi su tipa \texttt{AwareDatetime}, pa zahtjev bez oznake zone ne prolazi.

\end{frame}

% ------------------------------------------------------------
\begin{frame}[fragile]{Provjera nad cijelim modelom}
Validator uz polje vidi samo polja provjerena prije njega, dakle ona deklarirana iznad. Usporedba koja se na taj redoslijed ne može osloniti piše se drukčije:

\begin{lstlisting}[language=python]
# app/schemas/competition.py
@model_validator(mode="after")
def deadlines_before_competition(self):
    if self.final_deadline.date() >= self.date:
        raise ValueError(<@"Rokovi moraju isteći prije dana natjecanja"@>)
    return self
\end{lstlisting}

\vspace{0.4em}
\texttt{model\_validator} s oznakom \texttt{mode="after"} izvodi se kad su sva polja popunjena, pa im pristupa preko \texttt{self}.

\vspace{0.9em}
Rok i dan natjecanja mogu stajati u shemi bilo kojim redom, a provjera i dalje radi.
\end{frame}

% ------------------------------------------------------------
\begin{frame}{Pravila koja traže bazu}
Shema ne može pogledati u bazu. Validator je obična funkcija bez sesije, a u asinkronom radu unutar nje nema ni čekanja odgovora baze.

\vspace{1.1em}
\textbf{Shema prosuđuje vrijednost samu za sebe, ne u odnosu na druge zapise.}

\vspace{1.1em}
Je li natjecatelj već prijavljen na to natjecanje, bez baze se ne zna. Takvo je pravilo posao sloja \codefile{services}, koji dolazi iza sheme.
\end{frame}

% ------------------------------------------------------------
\begin{frame}[fragile]{Izmjena i neposlana polja (I)}
\begin{lstlisting}[language=python]
# app/schemas/lifter.py
class LifterUpdate(BaseModel):
    first_name: str | None = Field(default=None, min_length=1, max_length=80)
    last_name: str | None = Field(default=None, min_length=1, max_length=80)
    birth_date: date | None = None
    gender: Literal["M", "F"] | None = None
\end{lstlisting}

\vspace{0.4em}
Sva su polja neobavezna, pa klijent šalje ona koja mijenja. Time nastaje dvoznačnost: izostalo polje i polje poslano kao \texttt{null} u shemi izgledaju jednako, a znače različito.
\end{frame}

% ------------------------------------------------------------
\begin{frame}[fragile]{Izmjena i neposlana polja (II)}
Provjerena shema stiže u servis kao \texttt{body}. Ondje se pretvara u rječnik i polje po polje prepisuje u zapis:

\begin{lstlisting}[language=python]
# app/services/lifter_service.py
async def update_lifter(db, club_id, lifter_id, body: LifterUpdate, user):
    lifter = await get_lifter(db, club_id, lifter_id, user)
    for field, value in body.model_dump(exclude_unset=True).items():
        setattr(lifter, field, value)
    await db.flush()
    return lifter
\end{lstlisting}

\vspace{0.4em}
\texttt{exclude\_unset} razdvaja ta dva slučaja: u rječnik ulaze samo polja koja su doista poslana. Bez te oznake ušla bi i neposlana, svako sa svojom zadanom vrijednošću \texttt{None}, pa bi izmjena jednoga polja izbrisala sva ostala.
\end{frame}

% ------------------------------------------------------------
\begin{frame}[fragile]{Izmjena i neposlana polja (III)}
Poslano kao \texttt{null} \texttt{exclude\_unset} ne odvaja: takvo polje ulazi u rječnik i upisuje \texttt{null} u stupac koji ga ne dopušta. Zato \texttt{LifterUpdate} dobiva još jedan redak:

\begin{lstlisting}[language=python]
# app/schemas/lifter.py
class LifterUpdate(BaseModel):
    first_name: str | None = Field(default=None, min_length=1, max_length=80)
    ...
    _not_null = field_validator(
        "first_name", "last_name", "birth_date", "gender", mode="before"
    )(reject_null)
\end{lstlisting}

\vspace{0.4em}
\texttt{reject\_null} je obična funkcija iz \codefile{app/schemas/common.py}: podigne \texttt{ValueError} ako je vrijednost \texttt{None}. Pydantic je pozove za svako od četiriju navedenih polja, i to samo ako je poslano.
\end{frame}

% ------------------------------------------------------------
\begin{frame}{Vježba: što čuvaju \texttt{exclude\_unset} i \texttt{reject\_null}}
Dvije izmjene, svaka po jedan redak.

\vspace{0.9em}
\begin{enumerate}
  \itemsep0.5em
  \item \texttt{PATCH} natjecatelja samo s poljem \texttt{first\_name}, pa \texttt{GET} istoga zapisa. Provjerite jesu li se ostala polja promijenila.
  \item U \codefile{lifter\_service} maknite \texttt{exclude\_unset=True} i ponovite isti \texttt{PATCH}, pa opet \texttt{GET}.
  \item Vratite \texttt{exclude\_unset=True}, pa zakomentirajte \texttt{\_not\_null} u \texttt{LifterUpdate} i pošaljite \texttt{"last\_name": null}.
  \item Vratite oboje.
\end{enumerate}

\vspace{0.9em}
Usporedite odgovore iz drugoga i trećega koraka s onim iz prvoga.
\end{frame}

% ------------------------------------------------------------
\begin{frame}[fragile]{Odgovor o prijavi}
\begin{lstlisting}[language=python]
# app/schemas/registration.py
class RegistrationResponse(BaseModel):
    id: int
    lifter_id: int
    competition_id: int
    category: str
    total: int
    status: str
    registered_at: datetime
    lifter: LifterBrief

    model_config = {"from_attributes": True}
\end{lstlisting}

\vspace{0.4em}
Uz vlastita polja prijave zadnje polje nosi i sam povezani zapis. Da odgovor sadrži samo \texttt{lifter\_id}, klijent bi za svako ime slao novi zahtjev.
\end{frame}

% ------------------------------------------------------------
\begin{frame}[fragile]{Suženi povezani zapis}
Shema natjecatelja, ali samo s onim što popis prijava treba:

\begin{lstlisting}[language=python]
# app/schemas/registration.py
class LifterBrief(BaseModel):
    first_name: str
    last_name: str

    model_config = {"from_attributes": True}
\end{lstlisting}

\vspace{0.4em}
Datum rođenja i klub u popisu prijava nikome ne trebaju. Shema odgovora je mjesto na kojem se odlučuje koja polja povezanog zapisa izlaze van, a koja ostaju u aplikaciji.

\vspace{0.7em}
\texttt{LifterBrief} je time četvrta shema natjecatelja. Deklarirana je uz prijavu, u \codefile{app/schemas/registration.py}, jer ondje služi.
\end{frame}

% ============================================================
\sekcija{Popisi}{Zašto popis ne vraća sve, straničenje i filtriranje, redoslijed, povezani zapisi u popisu.}

% ------------------------------------------------------------
\begin{frame}{Popis bez granice}
Zahtjev \texttt{GET /clubs/12/lifters} ne kaže koliko zapisa traži. Dok ih klub ima desetak, odgovor je nekoliko kilobajta i nitko ne primijeti ništa.

\vspace{1.1em}
Kod kluba s tri tisuće natjecatelja isti zahtjev ima tri posljedice:

\vspace{0.7em}
\begin{itemize}
  \itemsep0.6em
  \item upit koji čita cijelu tablicu
  \item odgovor koji cijel putuje mrežom
  \item klijent koji prikaže prvih dvadeset, a ostalo odbaci
\end{itemize}

\vspace{1.1em}
Na razvojnoj bazi s troje natjecatelja ništa od toga nije vidljivo.
\end{frame}

% ------------------------------------------------------------
\begin{frame}{Straničenje}
Odgovor se zato vraća u dijelovima. Klijent kaže koliko zapisa želi i koliko ih preskače:

\vspace{1.1em}
\begin{center}
\begin{tikzpicture}[
  cel/.style={draw=lineclr, thick, fill=white, minimum width=7mm,
              minimum height=7mm, font=\scriptsize, text=mutedclr},
  okvir/.style={draw=ink, very thick, rounded corners},
  opis/.style={font=\scriptsize, text=textclr}
]
  \foreach \i in {1,...,9} {
    \node[cel] (c\i) at (\i*7.6mm, 0) {\i};
  }
  \draw[okvir] ([xshift=-1.2mm,yshift=1.2mm]c1.north west)
    rectangle ([xshift=1.2mm,yshift=-1.2mm]c3.south east);
  \draw[okvir, draw=ink2, dashed] ([xshift=-1.2mm,yshift=1.2mm]c4.north west)
    rectangle ([xshift=1.2mm,yshift=-1.2mm]c6.south east);
  \node[opis, above=4mm of c2] {prva stranica};
  \node[opis, above=4mm of c5] {druga stranica};
  \node[opis, below=4mm of c2] {\texttt{offset=0}};
  \node[opis, below=4mm of c5] {\texttt{offset=3}};
\end{tikzpicture}
\end{center}

\vspace{1.1em}
Obje stranice imaju \texttt{limit=3}, a razlikuju se po \texttt{offset}-u: \texttt{limit} je širina prozora, \texttt{offset} koliko se zapisa preskoči prije njega. Vraćanje popisa u dijelovima naziva se straničenje (engl. \emph{pagination}).
\end{frame}

% ------------------------------------------------------------
\begin{frame}[fragile]{Parametri upita}
Granicu i pomak klijent šalje u adresi, iza upitnika:

\begin{lstlisting}
GET /clubs/1/lifters?limit=3&offset=0
GET /clubs/1/lifters?limit=3&offset=3
\end{lstlisting}

\vspace{0.6em}
Svaki par \texttt{ime=vrijednost} iza \texttt{?} jedan je parametar upita, a međusobno se odvajaju znakom \texttt{\&}.

\vspace{0.9em}
Parametri nisu dio putanje: obje adrese vode na istu rutu, samo s različitim vrijednostima. Zato ih ruta ne prepoznaje po mjestu nego po imenu.
\end{frame}

% ------------------------------------------------------------
\begin{frame}[fragile]{Granica u potpisu rute}
\begin{lstlisting}[language=python]
# app/routers/lifters.py
@router.get("/", response_model=list[LifterResponse])
async def list_lifters(
    club_id: int,
    db: AsyncSession = Depends(get_db),
    user: User = Depends(require_role("admin", "club")),
    limit: int = Query(default=20, ge=1, le=100),
    offset: int = Query(default=0, ge=0),
    ...
):
    ...
\end{lstlisting}

\vspace{0.3em}
\texttt{Query} opisuje parametar upita jednako kao što shema opisuje polje tijela: zadana vrijednost, najmanja i najveća dopuštena.
\end{frame}

% ------------------------------------------------------------
\begin{frame}{Zadane vrijednosti i gornja granica}
\begin{itemize}
  \item \textbf{Klijent koji o straničenju ne zna ništa dobio bi neograničen odgovor.} Zato obje vrijednosti imaju zadanu, pa se ruta ne može ni slučajno pozvati bez granice.
  \item \textbf{Klijent može zatražiti milijun zapisa.} Zato \texttt{limit} ima gornju granicu \texttt{le=100}, inače se ruta vraća na ponašanje s početka.
  \item \textbf{Klijent može zatražiti nula zapisa ili negativan pomak.} Zato stoje \texttt{ge=1} i \texttt{ge=0}, pa framework takav zahtjev odbija prije nego funkcija počne raditi.
\end{itemize}

\vspace{0.4em}
\begin{konvencija}
Nazivi \texttt{limit} i \texttt{offset} nisu propisani. U drugim se API-jima sreću i \texttt{page} i \texttt{per\_page}, a neki uz popis vraćaju i ukupan broj zapisa. Izbor je slobodan, ali mora biti isti na svim rutama istoga API-ja.
\end{konvencija}
\end{frame}

% ------------------------------------------------------------
\begin{frame}[fragile]{Filtriranje popisa}
Klub koji traži samo natjecateljice bez filtra dohvaća sve pa polovicu odbacuje. Zato ruta prima i uvjet pretrage:

\begin{lstlisting}[language=python]
# app/routers/lifters.py
    gender: str | None = Query(default=None, pattern=r"^[MF]$"),
\end{lstlisting}

\vspace{0.6em}
Zadana vrijednost \texttt{None} znači da je filtar neobavezan: izostane li, ruta vraća natjecatelje obaju spolova.

\vspace{0.9em}
\texttt{pattern} propušta samo \texttt{"M"} i \texttt{"F"}. Bez njega bi \texttt{?gender=Z} prošao do upita, koji bi uredno vratio prazan popis, pa bi se pogrešno napisan zahtjev činio ispravnim.
\end{frame}

% ------------------------------------------------------------
\begin{frame}[fragile]{Filtar i granica u upitu}
\begin{lstlisting}[language=python]
# app/repositories/lifter_repo.py
stmt = select(Lifter).where(Lifter.club_id == club_id)
if gender is not None:
    stmt = stmt.where(Lifter.gender == gender)
stmt = stmt.order_by(Lifter.id).limit(limit).offset(offset)
\end{lstlisting}

\vspace{0.3em}
Upit se gradi postupno: pripadnost klubu provjerava se uvijek, a spol samo ako je filtar poslan. Ruta i servis prosljeđuju vrijednosti, a \texttt{select} i nazivi stupaca ne izlaze iz sloja \codefile{repositories}.

\vspace{0.3em}
Filtar i granica tako djeluju na upit, a ne na već dohvaćeni popis.
\end{frame}

% ------------------------------------------------------------
\begin{frame}{Poredak kao uvjet straničenja}
Straničenje ne može ispravno raditi bez izrijekom zadanog poretka zapisa, jer ga baza sama ne jamči.

\vspace{1.1em}
Bez \texttt{order\_by} postoji mogućnost da isti zapis dođe i na prvoj i na drugoj stranici, dok drugi ne dođe ni na jednu. Klijent to ne može primijetiti, jer svaka stranica izgleda uredno.

\vspace{1.1em}
Koliko redaka baza mora pročitati da vrati stotu stranicu?

\pause
\begin{konvencija}
Uz \texttt{offset} baza pri svakom zahtjevu pročita i odbaci sve prethodne retke, pa duboke stranice postaju spore. Drugi pristup pamti oznaku zadnjeg viđenog zapisa i traži one iza nje: brži je i otporan na umetanje, ali ne dopušta skok na proizvoljnu stranicu. U našem se projektu koristi \texttt{offset}, jer su popisi kratki.
\end{konvencija}
\end{frame}

% ------------------------------------------------------------
\begin{frame}[fragile]{Odgovor koji traži dvije tablice}
U bazi prijava ne drži natjecatelja nego samo njegovu oznaku:

\begin{lstlisting}[language=python]
# app/models/registration.py
lifter_id: Mapped[int] = mapped_column(ForeignKey("lifters.id"))
\end{lstlisting}

\vspace{0.5em}
Shema odgovora traži i sam zapis:

\begin{lstlisting}[language=python]
# app/schemas/registration.py
class RegistrationResponse(BaseModel):
    ...
    lifter_id: int
    lifter: LifterBrief
\end{lstlisting}

\vspace{0.5em}
Oznaka nije zapis. Da bi od broja \texttt{31} nastao \texttt{LifterBrief}, netko mora pročitati i tablicu \texttt{lifters}.
\end{frame}

% ------------------------------------------------------------
\begin{frame}[fragile]{Lijeno učitavanje povezanog zapisa}
Taj posao ORM nudi sam. U modelu uz strani ključ stoji i veza:

\begin{lstlisting}[language=python]
# app/models/registration.py
lifter: Mapped[Lifter] = relationship(back_populates="registrations")
competition: Mapped[Competition] = relationship(
    back_populates="registrations"
)
\end{lstlisting}

\vspace{0.2em}
Ta dva retka nisu stupci tablice nego veze. Zahvaljujući njima povezani se zapis čita kao svojstvo, dakle \texttt{registration.lifter}, bez pisanja upita.

\vspace{0.2em}
ORM taj zapis ne dohvaća unaprijed nego tek u trenutku čitanja svojstva. Taj se postupak naziva lijeno učitavanje (engl. \emph{lazy loading}).
\end{frame}

% ------------------------------------------------------------
\begin{frame}[fragile]{Lijeno učitavanje u asinkronom radu}
U sinkronom radu lijeno učitavanje prolazi neopaženo: dretva pošalje upit, pričeka odgovor i nastavi. Čekanje je dio samog čitanja svojstva i nigdje se ne vidi.

\vspace{0.5em}
U asinkronom radu svako čekanje baze mora biti napisano:

\begin{lstlisting}[language=python]
reg = await registration_repo.get_by_id(db, 7)  <@\kom{čeka bazu}@>
ime = reg.lifter.first_name                     <@\kom{ne čeka ništa}@>
\end{lstlisting}

\vspace{0.4em}
Prvi je redak poziv funkcije, pa uz njega \texttt{await} može stajati. Drugi je čitanje atributa, koje Python izvede odmah: ORM bi upit morao poslati unutar te točke, a ondje se na bazu više nema kako čekati.

\vspace{0.5em}
Zato \texttt{reg.lifter} ne vrati podatke nego podigne iznimku, i to tek kad se shema odgovora pokušava popuniti, pa poruka upućuje na mjesto daleko od uzroka.
\end{frame}

% ------------------------------------------------------------
\begin{frame}[fragile]{Dohvat povezanih zapisa unaprijed}
U sinkronom radu lijeno učitavanje ne bi puklo nego poskupjelo: popis od pedeset prijava izveo bi pedeset i jedan upit, jedan za popis i po jedan za natjecatelja. To se naziva problemom N+1.

\vspace{0.3em}
Asinkroni rad tu mogućnost uklanja: povezane zapise treba zatražiti u istom dohvatu, dok se na bazu još smije čekati:

\begin{lstlisting}[language=python]
# app/repositories/registration_repo.py
stmt = (
    select(Registration)
    .options(selectinload(Registration.lifter))
    .where(Registration.competition_id == comp_id)
)
\end{lstlisting}

\vspace{0.3em}
\texttt{selectinload} nalaže ORM-u da nakon upita za prijave izvede još jedan, kojim odjednom dohvaća sve natjecatelje. Dva upita ukupno, bez obzira na broj prijava.
\end{frame}

% ------------------------------------------------------------
\begin{frame}[fragile]{Vježba: priprema prijava}
Prazan popis nema što pokazati, pa ga prvo treba napuniti. Prijavite se kao klub \texttt{behemot@plreg.hr}.

\vspace{0.8em}
\begin{enumerate}
  \itemsep0.5em
  \item \texttt{GET /competitions} daje oznaku natjecanja \emph{Zimski kup}.
  \item \texttt{GET /clubs} daje oznaku vašega kluba, a \texttt{GET /clubs/\{id\}/lifters} oznake Ane i Ivana.
  \item Oboje prijavite na natjecanje, svakoga svojim zahtjevom:
\end{enumerate}

\begin{lstlisting}
POST /competitions/<oznaka natjecanja>/registrations

{"lifter_id": <oznaka Ivana>, "category": "83", "total": 0}
{"lifter_id": <oznaka Ane>,   "category": "63", "total": 0}
\end{lstlisting}

\vspace{0.4em}
Kategorija mora odgovarati spolu natjecatelja, inače odgovor je \texttt{400}.
\end{frame}

% ------------------------------------------------------------
\begin{frame}{Vježba: brojanje upita}
\begin{enumerate}
  \itemsep0.5em
  \item U \codefile{app/core/database.py} dodajte \texttt{echo=True} u \texttt{create\_async\_engine} i ponovno pokrenite poslužitelj.
  \item Dohvatite popis prijava i u konzoli prebrojite \texttt{SELECT} naredbe.
  \item Pronađite među njima upit nad tablicom \texttt{lifters} i pročitajte što stoji u njegovu \texttt{WHERE ... IN (...)}.
  \item Vratite \texttt{echo} na staro.
\end{enumerate}

\vspace{1.1em}
Koliko bi upita nad \texttt{lifters} bilo da natjecanje ima pedeset prijava?
\end{frame}

% ------------------------------------------------------------
\begin{frame}{Sažetak}

\begin{itemize}
  \item Resurs je vrsta zapisa objavljena pod vlastitom adresom; preslikavanje iz tablice ne mora biti jedan na jedan.
  \item Nad svakim se resursom ponavljaju iste četiri operacije; stvaranje vraća zapis jer oznaku dodjeljuje baza, brisanje ne vraća ništa.
  \item Ugniježđena adresa nosi pripadnost; kod samostalne je pripadnost uvjet pretrage koji se mora provjeriti.
  \item Unos, izmjena i odgovor tri su različite sheme, a \texttt{exclude\_unset} razlikuje neposlano polje od poslanog praznog.
  \item Shema prosuđuje vrijednost samu za sebe; pravilo koje traži uvid u bazu pripada sloju iznad nje.
  \item Popis se vraća u dijelovima, s gornjom granicom i zadanim redoslijedom; povezani zapisi traže se unaprijed.
\end{itemize}
\end{frame}

\end{document}
